\section{Erd\H{o}s Problem \#32: independent continuation}
\noindent Date: 23 September 2026. Source versions read in full: \texttt{32.tex} in \texttt{gpt\_pro\_5.4}.

\subsection*{1) OBJECTIVE AND SOURCE AUDIT}
The target is a cofinite additive complement $A$ to the primes with
$A(x)=o((\log x)^2)$, ideally $O(\log x)$. The supplied file proves a
different, two-thin-summand conclusion using an external theorem. Its
parity-shift step is valid, but replacing $A$ by $T+T$ restores the quadratic
logarithmic bound. I do not count that as closing the original gap, and do
not use its unverified random-model assertion as a theorem about primes.

\subsection*{2) WORK: A MULTIPLICITY OBSTRUCTION}
Here is a precise condition that any near-minimal complement must evade.
Let $A\subset\mathbb N_0$ be fixed and
\[
 r(n)=\#\{a\in A:n-a\text{ is prime}\},\qquad
 C(x)=\sum_{n\le x}\binom{r(n)}2.
\]
Suppose $r(n)\ge1$ for $n\ge n_0$, and suppose also that $r(n)\le K$
for these $n$, for a fixed integer $K\ge1$. Then
\[
 A(x)\pi(x)\ge x-O(1)+\frac{2C(x)}K-O(1).
\]
Indeed, interchange the order of summation to obtain
$\sum_{n\le x}r(n)=\sum_{a\le x}\pi(x-a)\le A(x)\pi(x)$.
For $1\le r\le K$,
$\binom r2\le K(r-1)/2$. Summing this inequality over $n_0\le n\le x$
proves the assertion; the finitely many earlier values contribute constants.
Consequently, if $C(x)\ge(\delta+o(1))x$ along a subsequence, then the
prime number theorem gives
\[
 A(x)\ge\left(1+\frac{2\delta}{K}+o(1)\right)\log x
\]
along that subsequence. The unconditional elementary first-order bound is
$A(x)\ge(1-o(1))\log x$, obtained without the multiplicity assumption.

The collision term has the exact form
\[
 C(x)=\sum_{\substack{a<b\\a,b\in A\cap[0,x]}}
 \#\{n\le x:n-a\text{ and }n-b\text{ are prime}\}.
\]
Thus an attempted construction close to the counting lower bound must
control actual shifted-prime correlations. Cardinality estimates alone
do not control this sum.

\subsection*{3) ADVERSARIAL VERIFICATION}
The bounded multiplicity hypothesis is additional and has not been proved
for an arbitrary thin complement. Nor have I proved a positive linear lower
bound for $C(x)$. The displayed obstruction is therefore conditional on
these two explicit assumptions, whereas the logarithmic lower bound is
unconditional. It is not a disproof of an $O(\log x)$ complement.

\subsection*{7) FINAL AND REFERENCES}
\textbf{UNRESOLVED.} This attempt provides a concrete overlap calculation
for the original one-shift problem, but neither a sufficiently thin set nor
a general obstruction to one. The supplied two-summand result remains a
different problem. Literature check: Dai--Pan, \emph{Note on the additive
complements of primes}, \url{https://arxiv.org/abs/1101.1653}; the search
record was checked, but no stronger theorem from it is used here.
The direct official problem-page fetch was unavailable.

\subsection*{8) COMPLETION ESTIMATE}
\textbf{COMPLETION: 8\%}
This is a subjective estimate of progress toward the entire stated problem, not a probability that the conjecture or an unverified proof is correct.
